A custom solver is implemented in C++ using the SCIP 10 framework with SoPlex 8~\cite{hojny_scip_2025}.
Note that the implementation of~\cite{lawless_interpretable_2023} is available on GitHub~\cite{lawless_conlawfaircg_2026}.
That implementation is writen in Python and uses proprietary software Gurobi.
Our implementation is fully open source and also available on \href{https://github.com/kalkoen/rule-sets-for-tf}{GitHub}.
This custom solver includes an efficient heuristic pricer and was designed to support the extensions discussed in the next chapters.
The custom solver has two modes:
\begin{enumerate}
    \item \textbf{Full:} Constructs and solves the full master problem \eqref{eq:BM}, with all possible rules (columns),
    \item \textbf{Column generation (CG):} Performs the price-then-branch approach.
\end{enumerate}
The Full mode can only realistically be solved for datasets with small $\J$ or small $D$.
It can serve to verify, to some extent, the performance and validity of the column generation approach.
The CG mode can use up to three pricers:
\begin{enumerate}
    \item \textbf{Iterative:} Iterates over all possible rules and computes reduced cost naively,
    \item \textbf{Exact:} Solves MIP \eqref{eq:BP} to optimality,
    \item \textbf{Heuristic:} Employs a beam search.
\end{enumerate}
The Exact pricer produces a single rule minimizing reduced cost (unless exact pricing is interrupted early).
The other pricers can be configured to return up to some maximum of rules, and to return a ``diverse'' collection of rules.

We now discuss the Heuristic pricer and the efforts taken to implement it efficiently.
The Heuristic pricer builds a tree, in which every node is a rule and each layer represents the rules of a given size.

The root node (at depth $0$) is the empty rule $\emptyset$, covering all points.
The Heuristic pricer starts by constructing all children of the root node and computing the reduced cost.
Then the nodes in this layer are pruned: only the nodes with the lowest reduced cost are kept.
The beam width at this depth of the tree determines the number of nodes kept.
This process repeats: in iteration $\nu$, the children of each node at depth $\nu-1$ are constructed.
Next, the number of nodes at level $\nu$ are reduced to the beam width of layer $\nu$.

Computing a node's reduced cost naively costs $\Theta(nd)$ operations,
where $d$ is the depth of the node, which is the rule's size.
For large datasets, this can be slow.
We note that there exists overlap between the computations of parents and those of their descendents.
We may exploit this overlap.
Let $k_\text{parent}, k_\text{child}$ be the rules corresponding to some parent and child node, respectively.
Recall that the reduced cost
\begin{equation}
        \hat \rho_{k_\text{child}} =  \sum_{i \in \cN} \Ind(k_\text{child} \in \K_i)
    -\sum_{i \in \cP}  \mu_i  \Ind(k_\text{child} \in \K_i) +\lambda  c_{k_\text{child}}.  \label{eq:BP_redcost_node}
\end{equation}
Define $\I_k$ as the set of points covered by rule $k$, so
\begin{equation}
    \I_k \coloneqq \{i \in \I \fw k \in \K_i\} = \left\{i \in \I \fw \bigwedge_{j \in k} X_{ij}\right\}, \tquad k \in \K.\label{eq:boolean_I_k}
\end{equation}
Define $\P_k$ and $ \N_k$ analogously.
Observe that
\[k_\text{parent} \notin \K_i \implies k_\text{child} \notin \K_i, \tquad i \in \I.\]
In other words, a point that does not satisfy the parent rule can never satisfy the child rule,
because the child rule is more restrictive.
Then
\[
    \I_{k_\text{child}} \subseteq \I_{k_\text{parent}}.
\]
The analogous property holds for $\P_k$ and $\N_k$.
It follows that
\begin{equation}
    \begin{split}
        \hat \rho_{k_\text{child}} = \; &  |\cN_{k_\text{child}}|
    -\sum_{i \in \cP_{k_\text{child}}}  \mu_i +\lambda c_{k_\text{child}} \\
        = \; & |\cN_{k_\text{parent}}| - |\cN_{k_\text{parent}} \cap \N_{k_\text{child}}| \\
    & -\left(\sum_{i \in \cP_{k_\text{parent}}} \mu_i - \sum_{i \in \cP_{k_\text{parent}} \cap \P_{k_\text{child}}} \mu_i \right) \\
        & +\lambda (c_{k_\text{parent}}+1) \\
        = \; & \hat{\rho}_{k_\text{parent}} - |\cN_{k_\text{parent}} \cap \N_{k_\text{child}}|
        + \sum_{i \in \cP_{k_\text{parent}}\cap \P_{k_\text{child}}} \mu_i
        + \lambda.
    \end{split}
    \label{eq:BP_redcost_node_deriv}
\end{equation}
To use this equation, we may store $\P_k$, $\N_k$ and the reduced cost for each node.
In our implementation, the sets are stored as bitmaps, reducing memory footprint and allowing for efficient treatment of
\[\cP_{k_\text{parent}} \cap \P_{k_\text{child}} \quad \text{and} \quad \N_{k_\text{parent}} \cap \N_{k_\text{child}}.\]
The cost of this approach is memory.
For every node kept in the search tree, these bitmaps need to be stored.
To avoid unnecessary memory use, these bitmaps are computed after pruning,
meaning these bitmaps are never created for pruned nodes.

The heuristic pricer can also be configured to do a ``deep'' search.
In this case, pruning is not done per layer but per parent node.
Setting the ``beam width'' (which is no longer a beam width, but a parent's degree) to $m$ or higher,
this effectively means the heuristic pricer does enumeration of all rules up to size $D$, only through this memory-heavy but computationally-efficient approach.
Depending on the available memory, the number of points number of features and $D$, the iterative pricer may be preferred.

We compiled our implementation on Eindhoven University of Technology's High Performance Cluster (HPC), allowing us to
perform many experiments simultaneously.
The HPC is a shared computer.
Hence, we cannot guarantee experiments are always run under the same load conditions.
Furthermore, we did not specifically request the same exact hardware to be used for each experiment.
Therefore, running times are not to be compared between experiments.

To test whether our implementation performs similarly to the implementation used in \cite{lawless_interpretable_2023}
we perform a benchmark.
We test our implementation on datasets marked ``small'' in the paper.
Some of the datasets could not be traced back.
The others were downloaded from the UCI Machine Learning Repository~\cite{kelly_markelle_uci_nodate}.
We set $C$ to the mean value reported in the paper, which was reportedly the average
$C$ yielding the best test accuracy.
For continuous variables, binarization is performed with $9$ quantiles.
Categorical variables are one-hot encoded.
We employ the heuristic pricer with beam widths $(100,50,50,20)$, identical to the paper.
The exact pricer is used when the heuristic pricer does not find any rule with negative reduced cost.
The pricing stage is limited to 2 hours.
The branching stage is limited to 1 hour, but in practice this limit is never reached.
Like in the original paper, 10-fold stratified cross validation is used.

\begin{adjustwidth}{-2cm}{-2cm}
\begin{table}[h]
    \begin{tabular}{l  V{1.7cm} r@{ }l r@{ }l r@{ }l r@{ }l}
\toprule
\parbox{1.7cm}{ \textbf{ Dataset}} & \textbf{$C$} & \multicolumn{2}{c}{\parbox{1.7cm}{\centering \textbf{Solve time (s)}}} & \multicolumn{2}{c}{\parbox{1.7cm}{\centering \textbf{Test acc. (\%)}}} & \multicolumn{2}{c}{\parbox{1.7cm}{\centering \textbf{Test acc. (\%) in \cite{lawless_interpretable_2023}}}} & \multicolumn{2}{c}{\parbox{1.7cm}{\centering \textbf{Train acc. (\%)}}} \\
\midrule
ILPD~\cite{ramana_ilpd_2022} & 11 & \num{6529} & (\num{1191}) & \num{69.4} & (\num{3.0}) & \num{69.6} & (\num{1.2}) & \num{78.2} & (\num{0.3}) \\
WDBC~\cite{wolberg_william_breast_1993} & 14 & \num{6859} & (\num{519}) & \num{93.3} & (\num{3.8}) & \num{94.0} & (\num{1.2}) & \num{99.0} & (\num{0.2}) \\
banknote~\cite{lohweg_banknote_2012} & 25 & \num{105} & (\num{35}) & \num{99.3} & (\num{0.7}) & \num{99.1} & (\num{0.3}) & \num{99.8} & (\num{0.1}) \\
heart~\cite{janosi_heart_1989} & 12 & \num{649} & (\num{303}) & \num{72.7} & (\num{7.7}) & \num{78.9} & (\num{2.4}) & \num{87.6} & (\num{0.9}) \\
ionosphere~\cite{sigillito_ionosphere_1989} & 13 & \num{7056} & (\num{456}) & \num{88.3} & (\num{5.1}) & \num{90.0} & (\num{1.8}) & \num{95.8} & (\num{0.6}) \\
tic-tac-toe~\cite{aha_tic-tac-toe_1991} & 32 & \num{1689} & (\num{567}) & \num{100.0} & (\num{0.0}) & \num{100.0} & (\num{0.0}) & \num{100.0} & (\num{0.0}) \\
transfusion~\cite{yeh_blood_2008} & 6 & \num{14} & (\num{14}) & \num{79.3} & (\num{3.9}) & \num{77.9} & (\num{1.4}) & \num{79.4} & (\num{0.5}) \\
\bottomrule
\end{tabular}

    \caption{Results for the benchmark between our implementation and the implementation of \cite{lawless_conlawfaircg_2026}.
    Metrics are reported in ``[Mean] ([Standard deviation])'' format. ``Accuracy'' is abbreviated to ``Acc.''.}
    \label{tab:benchmark_accuracy}
\end{table}
\end{adjustwidth}

The results are given in \autoref{tab:benchmark_accuracy}.
Test accuracies are similar for both implementations.
Not evident from the table is that our implementation is slower.
In some cases, the time limit of the pricing stage is reached.
We attribute this to the speed at which the exact pricer solves the pricing problem.